
\textbf{Lemme 1.10.} \uline{Avec les notations et les hypothèses de} 1.1, \uline{soient} @@M0@@ et @@M1@@ \uline{deux faisceaux localement libres de type fini sur} @@M2@@, \uline{tels que} :
@@M3@@
\uline{alors il existe un entier} @@M4@@, \uline{tel que pour} @@M5@@, \uline{on ait} :
@@M6@@

Soit @@M7@@ le groupe abélien libre engendré par les classes (à isomorphisme près) de faisceaux localement libres de type fini sur @@M8@@. Dire que, dans @@M9@@, @@M10@@ et @@M11@@ ont même classe, c'est dire, d'après la construction du @@M12@@, qu'il existe des suites exactes de @@M13@@-Modules localement libres de type fini :
@@M14@@
et des entiers @@M15@@, tels que dans @@M16@@ on ait l'égalité :
@@M17@@
Pour tout @@M18@@, la tensorisation par @@M19@@ définit un automorphisme de @@M20@@, d'où :
@@M21@@
les suites @@M22@@ restant exactes.

Soient @@M23@@ le groupe abélien libre engendré par les classes (à isomorphisme près) de faisceaux quasi-cohérents sur @@M24@@, et @@M25@@ le sous-groupe engendré par les classes de faisceaux localement libres de type fini sur @@M26@@. Par image directe, on définit un homomorphisme @@M27@@, et on a donc encore :
